%!TEX TS-program = lualatex
%!TEX encoding = UTF-8 Unicode

\documentclass[11pt]{exam}

%\printanswers


\usepackage[left=1in,right=0.75in,top=1in]{geometry}     

\usepackage{fontspec}
\def\mainfont{Linux Libertine O}
\setmainfont[Ligatures=TeX, Contextuals={NoAlternate}, BoldFont={* Bold}, ItalicFont={* Italic}, Numbers={Proportional}]{\mainfont}
\setmonofont[Scale=MatchLowercase]{Inconsolata} 
\setsansfont[Scale=MatchLowercase]{Linux Biolinum O} 
\usepackage{microtype}

\usepackage{graphicx}
	\graphicspath{%
	{/Users/goby/Pictures/teach/478/exams/}}%

% Used to get the last two digits of the year for the header below.
\def\short#1{\csname @gobbletwo\expandafter\endcsname\number#1}

\pagestyle{headandfoot}
%\firstpageheader{\bfseries{Ichthyology, Exam 1, F14}}{}{\bfseries{Name: \enspace \makebox[2.5in]{\hrulefill}}}
\firstpageheader{\bfseries{BI 300: Hardy-Weinberg Quiz}}{}{%
	\ifprintanswers\textbf{KEY}\else\bfseries{Name: \enspace \makebox[2.5in]{\hrulefill}}\fi}
\runningheader{}{}{\small{pg. \thepage}}
%\runningheadrule

\footer{}{}{}%{\makebox[0.75in]{\hrulefill}}
\extrafootheight{-0.5in}

\usepackage{multicol}
\usepackage{booktabs}


%% Remove comment %% to print answer key.
\unframedsolutions
\renewcommand{\solutiontitle}{}
\SolutionEmphasis{\bfseries}

\pointsinmargin
%\marginpointname{ pts}
\marginpointname{ pt}


%% Create a Matching question format
\newcommand*\Matching[1]{
\ifprintanswers
	\textbf{#1}
\else
	\rule{2.1in}{0.5pt}
\fi
}
\newlength\matchlena
\newlength\matchlenb
\settowidth\matchlena{\rule{2.1in}{0pt}}
\newcommand\MatchQuestion[2]{%
	\setlength\matchlenb{\linewidth}
	\addtolength\matchlenb{-\matchlena}
	\parbox[t]{\matchlena}{\Matching{#1}}\enspace\parbox[t]{\matchlenb}{#2}}



\begin{document}

\begin{questions}

\question[3]
You should have reviewed the pages in your text on the Hardy-Weinberg Principle. Based on your review or prior knowledge, tell in your own words what is meant when a locus is said to be in Hardy-Weinberg Equilibrium (HWE).

\vspace{4cm}

\question[2]
The Hardy-Weinberg Principle is based on two simple, related mathematical formulas. Write them here. 

\vspace{1cm}

a. \vspace{1.5cm}

b. \vspace{1cm}

\question[5]
Tell what each mathematical term in each formula represents for a population in HWE.

\vspace{1cm}

$p$\vspace{1.5cm}

$q$\vspace{1.5cm}

$p^2$\vspace{1.5cm}

$2pq$\vspace{1.5cm}

$q^2$

\newpage
\question[4]
The Hardy-Weinberg Equations have five assumptions. Your text names all five but suggests that only four apply to the equations. List the four assumptions mentioned by your text. List the fifth for 1 point extra credit. You do not have to explain the assumptions, just list them.

\vspace{1cm}

a. \vspace{1.5cm}

b. \vspace{1.5cm}

c. \vspace{1.5cm}

d. \vspace{1.5cm}

E.C.

\vspace{1.5cm}

\end{questions}
\end{document}